\begin{theorem}[固定分辨率 Pareto 冷端的对数尖点]\label{thm:conclusion-fixed-resolution-pareto-coldend-log-cusp}
\leanverified{paper\_conclusion\_fixed\_resolution\_pareto\_coldend\_log\_cusp}
沿用定理 \ref{thm:pom-finite-pareto-legendre-curvature} 的记号，并设 \(d_m\) 非常值。记顶端三层多重度与对应基数为
\[
M_1:=M_m,\qquad
M_2:=M_m^{(2)},\qquad
M_3:=M_m^{(3)},
\]
\[
N_1:=\#\{x\in X_m:\ d_m(x)=M_1\},
\qquad
N_2:=\#\{x\in X_m:\ d_m(x)=M_2\},
\]
其中约定若 \(d_m\) 只有两层取值，则 \(M_3:=0\)。再记
\[
\kappa_\ast:=\log M_1,
\qquad
H_\ast:=\log N_1,
\qquad
\Delta_m:=\log\frac{M_1}{M_2}>0,
\qquad
a_m:=\frac{N_2}{N_1}.
\]
则当 \(\kappa\uparrow\kappa_\ast\) 时，固定分辨率 Pareto 上边界满足精确尖点展开
\[
\mathcal H_m(\kappa)
=
H_\ast
+
\frac{\kappa_\ast-\kappa}{\Delta_m}
\left(
\log\frac1{\kappa_\ast-\kappa}
+
1+\log(a_m\Delta_m)
\right)
+
o(\kappa_\ast-\kappa).
\]
因此冷端并非二次光滑闭合，而是一个由第二层高度差 \(\Delta_m\) 与第二层相对简并比 \(a_m\) 共同强制的对数尖点。
\end{theorem}
\begin{proof}
由定理 \ref{thm:pom-finite-pareto-legendre-curvature}，
\[
\kappa=F_m'(q),
\qquad
\mathcal H_m(\kappa)=F_m(q)-qF_m'(q),
\qquad
F_m(q):=\log S_q(m),
\]
且当 \(q\to+\infty\) 时有 \(\kappa\uparrow \kappa_\ast=\log M_1\)。

把幂和按顶端三层分拆，得到
\[
S_q(m)
=
N_1M_1^q+N_2M_2^q+O(M_3^q)
=
N_1M_1^q\Bigl(1+a_m\rho^q+O(\sigma^q)\Bigr),
\]
其中
\[
\rho:=\frac{M_2}{M_1}=e^{-\Delta_m}\in(0,1),
\qquad
\sigma:=\frac{M_3}{M_1}\in[0,\rho).
\]
于是
\[
F_m(q)
=
q\log M_1+\log N_1+\log\!\Bigl(1+a_m\rho^q+O(\sigma^q)\Bigr).
\]
对 \(q\) 求导得到
\[
\kappa_\ast-\kappa
=
-\frac{d}{dq}\log\!\Bigl(1+a_m\rho^q+O(\sigma^q)\Bigr)
=
a_m\Delta_m\,\rho^q+o(\rho^q).
\]
因此
\[
\rho^q=\frac{\kappa_\ast-\kappa}{a_m\Delta_m}+o(\kappa_\ast-\kappa),
\]
并且
\[
q
=
\frac1{\Delta_m}
\left(
\log(a_m\Delta_m)-\log(\kappa_\ast-\kappa)
\right)
+o(1).
\]

另一方面，由
\[
\mathcal H_m(\kappa)=F_m(q)-qF_m'(q)
\]
与上式展开可得
\[
\mathcal H_m(\kappa)
=
\log N_1+a_m(1+q\Delta_m)\rho^q+o(\rho^q).
\]
再将前述 \(\rho^q\) 与 \(q\) 的表达式代入并整理，即得
\[
\mathcal H_m(\kappa)
=
H_\ast
+
\frac{\kappa_\ast-\kappa}{\Delta_m}
\left(
\log\frac1{\kappa_\ast-\kappa}
+
1+\log(a_m\Delta_m)
\right)
+
o(\kappa_\ast-\kappa).
\]
证毕。
\end{proof}

\begin{corollary}[冷端有限曲率延拓不可能与二阶极点律]\label{cor:conclusion-pareto-coldend-curvature-pole}
在定理 \ref{thm:conclusion-fixed-resolution-pareto-coldend-log-cusp} 的设定下，当 \(\kappa\uparrow\kappa_\ast\) 时有
\[
\mathcal H_m'(\kappa)
=
-\frac1{\Delta_m}\log\frac1{\kappa_\ast-\kappa}
+O(1),
\]
以及
\[
\mathcal H_m''(\kappa)
=
-\frac{1}{\Delta_m(\kappa_\ast-\kappa)}
+o\!\left(\frac1{\kappa_\ast-\kappa}\right).
\]
特别地，\(\mathcal H_m\) 不可能在 \(\kappa_\ast\) 处具有有限曲率的 \(C^2\)-延拓；冷端几何必为真正的对数尖点。
\leanverified{paper_conclusion_pareto_coldend_curvature_pole_seeds}
\leanverified{paper_conclusion_pareto_coldend_curvature_pole_package}
\end{corollary}
\begin{proof}
由定理 \ref{thm:pom-finite-pareto-legendre-curvature}，
\[
\mathcal H_m'(\kappa)=-q,
\qquad
\mathcal H_m''(\kappa)=-\frac1{F_m''(q)}.
\]
而在定理 \ref{thm:conclusion-fixed-resolution-pareto-coldend-log-cusp} 的证明中已经得到
\[
\kappa_\ast-\kappa=a_m\Delta_m\rho^q+o(\rho^q),
\qquad
q=
\frac1{\Delta_m}\log\frac1{\kappa_\ast-\kappa}
+O(1).
\]
这立刻给出第一式。再由
\[
F_m''(q)=a_m\Delta_m^2\rho^q+o(\rho^q)
=
\Delta_m(\kappa_\ast-\kappa)+o(\kappa_\ast-\kappa),
\]
代入 \(\mathcal H_m''(\kappa)=-1/F_m''(q)\) 即得第二式。证毕。
\end{proof}

\begin{theorem}[冷端尖点对顶端第二层参数的层析]\label{thm:conclusion-pareto-coldend-two-parameter-tomography}
\leanverified{paper\_conclusion\_pareto\_coldend\_two\_parameter\_tomography}
在定理 \ref{thm:conclusion-fixed-resolution-pareto-coldend-log-cusp} 的设定下，冷端尖点满足
\[
\lim_{\kappa\uparrow\kappa_\ast}
\frac{\mathcal H_m(\kappa)-H_\ast}
{(\kappa_\ast-\kappa)\log\frac1{\kappa_\ast-\kappa}}
=
\frac1{\Delta_m},
\]
并且
\[
\lim_{\kappa\uparrow\kappa_\ast}
\left[
\frac{\Delta_m\bigl(\mathcal H_m(\kappa)-H_\ast\bigr)}{\kappa_\ast-\kappa}
-
\log\frac1{\kappa_\ast-\kappa}
\right]
=
1+\log(a_m\Delta_m).
\]
因此，冷端 Pareto 前沿的局部几何已经唯一恢复
\[
\Delta_m=\log\frac{M_1}{M_2},
\qquad
a_m=\frac{N_2}{N_1}.
\]
\end{theorem}
\begin{proof}
把定理 \ref{thm:conclusion-fixed-resolution-pareto-coldend-log-cusp} 的展开写成
\[
\mathcal H_m(\kappa)-H_\ast
=
\frac{y}{\Delta_m}
\left(
\log\frac1y+1+\log(a_m\Delta_m)
\right)
+o(y),
\qquad
y:=\kappa_\ast-\kappa.
\]
分别除以 \(y\log(1/y)\) 与 \(y\)，便得到两个极限。第一个极限给出 \(\Delta_m^{-1}\)，第二个极限则给出 \(a_m\Delta_m\)。故 \(\Delta_m\) 与 \(a_m\) 可被唯一恢复。证毕。
\end{proof}

\begin{theorem}[右端可见相的支撑函数闭包]\label{thm:conclusion-right-edge-visible-phases-support-function-closure}
\leanverified{paper\_conclusion\_right\_edge\_visible\_phases\_support\_function\_closure}
沿用定理 \ref{thm:pom-spectrum-right-edge-support-function} 的记号，并假设定理 \ref{thm:pom-multiwell-softmax-low-temperature-expansion} 中的右端可见相有限，记为
\[
(\alpha_i,h_i)_{i=1}^{K},
\qquad
\alpha_1=\alpha_\ast>\alpha_2>\cdots>\alpha_K,
\qquad
h_i=f(\alpha_i).
\]
定义冷端缺口函数
\[
R(a):=P_a-a\alpha_\ast
\qquad
(a\in\ZZ_{\ge 1}).
\]
则有精确闭式
\[
R(a)=\max_{1\le i\le K}\bigl\{h_i-a(\alpha_\ast-\alpha_i)\bigr\}.
\]
因此 \(R(a)\) 是一条凸、分段仿射、单调不增的右端支撑函数；其每一段仿射片恰对应一个可见相，而任意两相的交叉尺度由
\[
a_{i\leftrightarrow j}
=
\frac{h_i-h_j}{\alpha_i-\alpha_j}
\]
给出。
\end{theorem}
\begin{proof}
由定理 \ref{thm:pom-spectrum-right-edge-support-function}，
\[
R(a)=\sup_{\alpha\in\Sigma_f}\bigl(f(\alpha)-a(\alpha_\ast-\alpha)\bigr).
\]
而在有限可见相假设下，右端有效支撑恰由有限个点 \((\alpha_i,h_i)\) 给出，故上确界退化为有限条仿射函数
\[
a\longmapsto h_i-a(\alpha_\ast-\alpha_i)
\]
的最大值。其余结论皆为有限条直线取上包络的直接推论。证毕。
\end{proof}

\begin{corollary}[不可见相与整数矩阶的暴露失败]\label{cor:conclusion-nonexposed-phase-fails-integer-domination}
在定理 \ref{thm:conclusion-right-edge-visible-phases-support-function-closure} 的设定下，若某一可见相 \((\alpha_j,h_j)\) 不是点集
\[
\{(\alpha_i,h_i)\}_{i=1}^{K}
\]
上凸包的暴露顶点，则对每个整数 \(a\ge 1\) 都有严格不等式
\[
h_j-a(\alpha_\ast-\alpha_j)<R(a).
\]
因此该相虽然可以在低温 softmax 展开中以指数小项出现，却不可能成为任何整数矩阶的主导相。换言之，整数矩阶真正看见的只是右端谱凸包的整数可见暴露相。
\leanverified{paper\_conclusion\_nonexposed\_phase\_integer\_domination\_seeds}
\leanverified{paper\_conclusion\_nonexposed\_phase\_integer\_domination\_package}
\end{corollary}
\begin{proof}
由定理 \ref{thm:conclusion-right-edge-visible-phases-support-function-closure}，
\[
R(a)=\max_{1\le i\le K}\bigl\{h_i-a(\alpha_\ast-\alpha_i)\bigr\}.
\]
若 \((\alpha_j,h_j)\) 不是上凸包的暴露顶点，则其对应仿射函数
\[
L_j(a):=h_j-a(\alpha_\ast-\alpha_j)
\]
在任何 \(a\ge 1\) 上都不可能成为上包络的支撑线。于是对每个整数 \(a\ge 1\) 都有
\[
L_j(a)<\max_i L_i(a)=R(a).
\]
证毕。
\end{proof}

\begin{theorem}[冷端凸包与大阶差分窗口的 Newton--Prony 等价]\label{thm:conclusion-coldend-convex-hull-difference-window-prony-equivalence}
\leanverified{paper\_conclusion\_coldend\_convex\_hull\_difference\_window\_prony\_equivalence}
沿用定理 \ref{thm:conclusion-right-edge-visible-phases-support-function-closure} 的设定，并记
\[
\mathcal C_{\mathrm{cold}}
:=
\operatorname{conv}^{\uparrow}\{(\alpha_i,h_i)\}_{i=1}^{K}
\]
为右端可见相的上凸包。进一步假设 \(\mathcal C_{\mathrm{cold}}\) 的每个暴露顶点的法锥都与 \(\ZZ_{\ge1}\) 相交。则以下三类数据彼此唯一决定：
\begin{enumerate}
  \item 冷端上凸包 \(\mathcal C_{\mathrm{cold}}\)；
  \item 冷端支撑函数
  \[
  R(a)=P_a-a\alpha_\ast
  \qquad (a\in\ZZ_{\ge1});
  \]
  \item 充分大 \(q\) 处的一段有限差分窗口
  \[
  \bigl\{\Delta P(q+j)\bigr\}_{j=0}^{K_0},
  \qquad
  \Delta P(q):=P(q+1)-P(q),
  \]
  其中 \(K_0\) 仅依赖于可见相个数上界。
\end{enumerate}
换言之，冷端凸包几何与大 \(q\) 差分谱在信息论上完全等价；前者是几何表述，后者是可审计算法表述。
\end{theorem}
\begin{proof}
\textup{(i)\(\Rightarrow\)(ii)} 由定理 \ref{thm:conclusion-right-edge-visible-phases-support-function-closure} 立即得到。

\textup{(ii)\(\Rightarrow\)(i)} 时，\(R(a)\) 是有限条仿射函数的上包络。每个暴露片都可写成
\[
a\longmapsto h_i-a(\alpha_\ast-\alpha_i).
\]
由仿射片的斜率与截距即可恢复相应的 \((\alpha_i,h_i)\)。在附加的整数可见性假设下，每个暴露顶点都在某个整数 \(a\) 上被支撑函数看见，因此 \(R(a)\) 的整数数据唯一恢复整个上凸包 \(\mathcal C_{\mathrm{cold}}\)。

\textup{(iii)\(\Rightarrow\)(i)} 时，定理 \ref{thm:pom-multiwell-softmax-low-temperature-expansion} 给出
\[
P(q)-q\alpha_1-h_1
=
\log\!\left(
1+\sum_{i\ge2}e^{-(\alpha_1-\alpha_i)q+(h_i-h_1)}
\right)
+O(e^{-cq}),
\]
故差分序列 \(\Delta P(q)\) 具有有限指数和的可审计扰动结构。再由定理 \ref{thm:pom-multiwell-top-parameter-identifiability} 的 Prony--Hankel 识别，可由充分大的有限差分窗口恢复全部可见竞争相 \((\alpha_i,h_i)\)，因而恢复其上凸包。

\textup{(i)\(\Rightarrow\)(iii)} 则只需把 \(\mathcal C_{\mathrm{cold}}\) 所对应的暴露相数据代回同一多井低温展开，即得到差分窗口的有限指数模型。证毕。
\end{proof}

\begin{theorem}[冷端支撑函数的有限整阶平台]\label{thm:conclusion-coldend-support-function-finite-integer-plateau}
\leanverified{paper\_conclusion\_coldend\_support\_function\_finite\_integer\_plateau}
在定理 \ref{thm:conclusion-right-edge-visible-phases-support-function-closure} 的设定下，定义
\[
a_{\mathrm{plat}}
:=
\max_{i\ge2}\frac{h_i-h_1}{\alpha_1-\alpha_i},
\qquad
\alpha_1=\alpha_\ast,
\qquad
h_1=h_\ast.
\]
则对每个整数 \(a>a_{\mathrm{plat}}\) 都有精确平台
\[
R(a)=h_1=h_\ast,
\qquad
P_a=a\alpha_\ast+h_\ast.
\]
因此冷端支撑函数在有限整阶之后出现真正的水平平台，而不是仅仅渐近接近。若再施加严格基态间隙并调用定理 \ref{thm:fold-pressure-freezing-threshold}，则这一有限整阶平台正是实参数冻结枝
\[
P_t=t\alpha_\ast+h_\ast
\qquad (t>a_{\mathrm c})
\]
在整数温度轴上的离散影像。
\end{theorem}
\begin{proof}
由定理 \ref{thm:conclusion-right-edge-visible-phases-support-function-closure}，
\[
R(a)=\max\Bigl\{h_1,\ \max_{i\ge2}\bigl[h_i-a(\alpha_1-\alpha_i)\bigr]\Bigr\}.
\]
对每个 \(i\ge2\)，一旦
\[
a>\frac{h_i-h_1}{\alpha_1-\alpha_i},
\]
便有
\[
h_i-a(\alpha_1-\alpha_i)<h_1.
\]
因此当 \(a>a_{\mathrm{plat}}\) 时，所有次级相对应的仿射线都严格落在主相水平线 \(h_1\) 之下，从而
\[
R(a)=h_1.
\]
于是
\[
P_a=a\alpha_1+h_1=a\alpha_\ast+h_\ast.
\]
若进一步满足定理 \ref{thm:fold-pressure-freezing-threshold} 的严格基态间隙假设，则对一切实数 \(t>a_{\mathrm c}\) 都有
\[
P_t=t\alpha_\ast+h_\ast,
\]
故上述整数平台正是该冻结枝在整数温度轴上的离散影像。证毕。
\end{proof}
